\section{System and Threat Model}\label{sec:overview}

\subsection{System boundary and actors}\label{sec:overview:layout}

Constitutional ARI-RQGM governs the institution surrounding an autonomous
research executor. The executor produces typed research artefacts,
measurements, review records, and manuscript candidates inside a
checkpoint. The governance system observes those records and controls
which prompt, component, and utility-policy versions may serve. It does
not control the host operating system, model weights, or the human who can
change the program.

The actors divide into four groups. \emph{Research actors} propose,
evaluate, defend, and write. \emph{Judiciary actors} collect evidence,
prosecute, defend governance motions, and adjudicate them. \emph{Meta
actors} generate candidate prompts and policies or select replay cases.
\emph{Fixed actors} validate constitutional rules, apply registry
transitions, repair affected records, and verify final claims. An actor
may produce a hashed record only within its declared role. The current
artifact does not implement digital signatures or role-key management.
Producing a recommendation never implies authority to enact it.

\subsection{Threat model}\label{sec:overview:pipeline}

The primary adversary is reflexive institutional optimisation: an
evolvable role may seek a change that makes its own outputs look better
rather than making the research more reliable. Concrete threats include
score inflation, metric gaming, evidence omission, prompt injection,
collusion between proposer and evaluator, a meta agent requesting broader
tools than its incumbent, a judge ruling on its own prosecution, and a
retired prompt remaining influential through derived scores.

The model also includes accidental failures. A process may stop between
registry events; a model may emit malformed output; an anchor corpus may
be absent; a prompt candidate may fail its output schema; or stored
evidence may lack the raw axes needed for re-scoring. These failures must
not produce a half-applied institution.

The threat model excludes a malicious host administrator, arbitrary
modification of the fixed program, compromise of the underlying hash
function, rollback of both a checkpoint and its local pin to the same old
prefix, secret modification of model weights by the provider, and
irreversible external API, device, or production-data operations.
Changing the constitutional program is treated as a human amendment, not
an attack the running system can prevent.

\subsection{Authority and verification boundary}\label{sec:overview:verification}

Authority flows in one direction. Model-bearing roles may create
proposals, attacks, defences, judgments, and candidates. The deterministic
kernel converts none of these directly into registry state. Instead it
checks the record, the acting role, the requested edge, epoch timing,
capability subset, provenance, and audit continuity. A transition engine
is the sole registry writer, and its output is checked again before
commit.

Two asymmetric failure rules preserve useful work without allowing an
uncertain institutional change. A governance hook on the research path is
\emph{fail-open}: its failure is recorded and candidate generation or
local computation may continue only in a side-effect-free sandbox. A
state-changing governance operation is \emph{fail-closed}: if its evidence
or legality cannot be established, the incumbent institution remains in
service. External irreversible actions require a separate fixed gate that
fails closed; the present kernel does not enforce that environmental
condition.

The final manuscript verifier is part of the fixed boundary. The
manuscript forward-declares numeric claim identifiers, formulas, operand
locations, units, and tolerances. The verifier recomputes this closed
formula vocabulary from execution evidence. It does not infer arbitrary
claims from natural language or establish the truth of non-numeric claims
or execution records. Writer and reviewer roles cannot change the
verifier, its inputs, or its verdict.

\subsection{State, provenance, and replay}\label{sec:overview:checkpoint}

The checkpoint is the unit of institutional identity. Its append-only
events are authoritative; registry, epoch, and erasure snapshots are
derived views that can be reconstructed by replay. Every governed record
carries the epoch, serving component, prompt hash, and relevant policy
hash. Audit events form a predecessor-linked hash chain. Consequently a
reader can distinguish what happened, which institution produced it, and
which later transition changed its eligibility. This local chain detects
interior deletion or modification, and a surviving local boundary pin
detects a shorter log. It cannot detect rollback of both log and pin to
the same old valid prefix. There is no external signed head, witness
checkpoint, or monotonic counter, and the hash is not an author signature.

Epoch transactions use prepare, close, open, and commit events. Replay
ignores an incomplete suffix and deterministically re-executes the
boundary. Founding registration is similarly write-once: a committed
founding transaction is replayed, never repeated. Logical erasure adds
staleness and eligibility fields rather than deleting bytes, preserving
the evidence needed to audit the retirement itself.

Epoch identity is split before it is composed. A policy fingerprint binds
the serving institution and all resolved governance settings, including
thresholds and budgets. An execution fingerprint binds the declared model,
decoding, tools, environment, and data snapshot. Missing provider or
environment revisions are explicitly marked unresolved. The former supports
institutional comparisons; the latter exposes when an apparently unchanged
policy ran on a different declared execution substrate. Neither turns an
opaque model service into a reproducible binary.

\subsection{Design principles}\label{sec:overview:principles}

Five principles organise the architecture.

\begin{enumerate}
  \item \textbf{Fixed law, evolvable institution.} The constitution,
        transition topology, and final verification path are outside the
        set of runtime-mutable artefacts.
  \item \textbf{Frozen service epochs.} A serving version changes only at a
        boundary. T16 emergency quarantine forms an immediate boundary and
        opens a new epoch fingerprint, so its before/after nodes have
        distinct institutional identities.
  \item \textbf{Separation of recommendation and authority.} Proposers,
        judges, and meta agents emit records. Only the transition engine
        writes the registry, under kernel validation.
  \item \textbf{History without inherited influence.} Retirement never
        deletes evidence, but affected descendants become stale or are
        re-scored before they may influence future selection.
  \item \textbf{Claims proportional to evidence.} Structural guarantees
        are stated as guarantees; mechanism demonstrations are not
        presented as evidence of improved research quality.
\end{enumerate}
